% !TEX root = ../../main.tex
%
% Chapter: Z_g(k) — Verlinde SL_2 closed form and Kummer–Bernoulli leading term
%
% Verlinde SL_2 partition function for affine sl_2 conformal blocks on a
% genus-g Riemann surface at level k.  The leading coefficient is the
% Hurwitz--Bernoulli scalar
% (-1)^g 2^{g-1} B_{2g-2}/(2g-2)!, and the first Bernoulli-index
% Kummer witnesses occur at B_12 and B_16.
%
% Author: Raeez Lorgat.

\chapter{The Verlinde polynomial \texorpdfstring{$Z_g(k)$}{Z\_g(k)}, Bernoulli leading term, and Kummer congruences}
\label{ch:z-g-kummer-bernoulli}

\ClaimStatusProvedHere

\section*{Orientation}

Fix $G = \mathrm{SL}_2$ and a level $k \in \mathbb{Z}_{\ge 0}$. Set
\begin{equation}
n \;=\; k + 2 \;=\; k + h^{\vee}.
\end{equation}
The Verlinde partition function on a closed oriented genus-$g$ surface,
equivalently the dimension of conformal blocks for the affine Kac--Moody
algebra $\widehat{\mathfrak{sl}}_2$ at level $k$, is
\begin{equation}
\label{eq:zg-verlinde-raw}
Z_g(k) \;=\; \sum_{\lambda} \left( \dim V_\lambda \right)^{2-2g},
\end{equation}
where $\lambda$ runs over integrable weights at level $k$ (i.e.\ spins
$\ell/2$ with $0 \le \ell \le k$). Using the modular $S$-matrix
$S_{\ell m} = \sqrt{2/n}\,\sin(\pi (\ell+1)(m+1)/n)$ and the Verlinde formula,
this rewrites as a sum over nonzero $S$-entries; Equation~\eqref{eq:zg-verlinde-raw}
is the standard starting point.

The arithmetic rests on four facts.

\begin{itemize}
\item Theorem~\ref{thm:z-g-closed-form-polynomial} rewrites $Z_g(k)$ as a
finite cosecant sum and promotes the cosecant sum to a polynomial of degree
$2g - 2$ in $n^2$.
\item Theorem~\ref{thm:z-g-polynomial-form} refines the polynomial to
$P_g(n) = n^{g-1}(n^2-1)\,R_{g-2}(n^2)$, degree $3g - 3$ in $n$.
\item Theorem~\ref{thm:z-g-leading-coefficient-bernoulli} identifies the
leading coefficient of $R_{g-2}$ as
$(-1)^g2^{g-1}B_{2g-2}/(2g-2)!$ via Hurwitz's cosecant series.
\item Theorem~\ref{thm:z-g-kummer-congruence} derives Kummer-type
divisibility witnesses: the irregular primes $691 \mid Z_7$ leading and
$3617 \mid Z_9$ leading.
\end{itemize}

The engine is Hurwitz's 1890 cosecant expansion, the Witten~1991 symplectic
volume interpretation, and the Verlinde~1988 fusion formula. The three
routes meet on the same leading coefficient with the same normalisation.

\section{Cosecant reduction: Theorem~\ref{thm:z-g-closed-form-polynomial}}

\begin{theorem}[$Z_g(k)$ closed form]
\label{thm:z-g-closed-form-polynomial}
\ClaimStatusProvedHere
For $G = \mathrm{SL}_2$, level $k \ge 0$, genus $g \ge 2$, and $n = k+2$,
\begin{equation}
\label{eq:zg-csc}
Z_g(k) \;=\; \left(\tfrac{n}{2}\right)^{g-1} S_{g-1}(n), \qquad
S_{g-1}(n) \;=\; \sum_{j=1}^{n-1} \csc^{2(g-1)}\!\!\left(\tfrac{\pi j}{n}\right).
\end{equation}
The function $S_{g-1}(n)$ is a polynomial of degree $2g-2$ in $n^2$ with
rational coefficients; equivalently $Z_g(k)$ is a polynomial of degree
$3g-3$ in $n$.
\end{theorem}

\begin{proof}
Step 1 (Verlinde). The modular $S$-matrix for $\widehat{\mathfrak{sl}}_2$ at
level $k$ has entries $S_{\ell m} = \sqrt{2/n}\,\sin(\pi(\ell+1)(m+1)/n)$ for
$0 \le \ell, m \le k$. The Verlinde formula reads
\begin{equation}
Z_g(k) \;=\; \sum_{\ell=0}^{k} S_{0\ell}^{2-2g}.
\end{equation}
Substitute $S_{0\ell} = \sqrt{2/n}\,\sin(\pi(\ell+1)/n)$ and relabel
$j = \ell+1 \in \{1, \ldots, k+1\} = \{1, \ldots, n-1\}$:
\begin{equation}
Z_g(k) \;=\; \sum_{j=1}^{n-1} \left(\tfrac{2}{n}\right)^{(2-2g)/2}\,
\sin^{2-2g}\!\left(\tfrac{\pi j}{n}\right)
\;=\; \left(\tfrac{n}{2}\right)^{g-1}\sum_{j=1}^{n-1}\csc^{2(g-1)}\!\!\left(\tfrac{\pi j}{n}\right).
\end{equation}
This yields~\eqref{eq:zg-csc}. \checkmark

Step 2 (polynomiality of $S_{g-1}(n)$ in $n^2$). The sum
$S_{g-1}(n) = \sum_{j=1}^{n-1} \csc^{2(g-1)}(\pi j / n)$ is a trigonometric
power sum over nontrivial $n$-th roots of unity. Writing
$\csc^2(\pi j/n) = -4\,\zeta^j / (1 - \zeta^j)^2$ with
$\zeta = e^{2\pi i/n}$, $S_{g-1}(n)$ is a symmetric function in
$\{\zeta, \zeta^2, \ldots, \zeta^{n-1}\}$. Classical results on sums of
cotangent/cosecant powers (Eisenstein; Cauchy) give that for each integer
$r \ge 1$,
\begin{equation}
\sum_{j=1}^{n-1} \csc^{2r}\!\!\left(\tfrac{\pi j}{n}\right)
\;=\; T_r(n^2),
\end{equation}
where $T_r \in \mathbb{Q}[x]$ is a polynomial of degree $r$ in its argument.
We record the first three:
\begin{align}
T_1(n^2) &= \tfrac{1}{3}(n^2 - 1), \label{eq:T1} \\
T_2(n^2) &= \tfrac{1}{45}(n^2-1)(n^2+11), \label{eq:T2} \\
T_3(n^2) &= \tfrac{1}{945}(n^2-1)(2 n^4 + 23 n^2 + 191). \label{eq:T3}
\end{align}
The only input used here is the classical cosecant-polynomiality
statement; the Bernoulli leading coefficient is computed independently
later from the Hurwitz expansion and is not needed for this
polynomiality step.

Taking $r = g-1$ gives $S_{g-1}(n) = T_{g-1}(n^2)$, a polynomial of degree
$g - 1$ in $n^2$, hence degree $2g-2$ in $n$. \checkmark

Step 3 (polynomiality of $Z_g$). Multiplying by $(n/2)^{g-1}$:
\begin{equation}
Z_g(k) \;=\; 2^{-(g-1)} n^{g-1} T_{g-1}(n^2),
\end{equation}
a polynomial in $n$ of degree $(g-1) + 2(g-1) = 3g-3$. \checkmark
\end{proof}

\begin{remark}[Witten normalisation]
\label{rem:witten-normalisation}
Witten~1991 identifies $Z_g(k)$ (times a level-independent factor) with the
symplectic volume of the moduli space of flat $\mathrm{SU}(2)$-connections
on $\Sigma_g$ in the large-$k$ limit. Specifically,
$\lim_{k\to\infty} Z_g(k) / n^{3g-3} = \mathrm{vol}(\mathcal{M}_g^{\mathrm{flat}})$
up to explicit Weyl-group factors. This pins the leading $n^{3g-3}$ term
against the symplectic-volume side and provides an independent verification of
the degree claim in Thm~\ref{thm:z-g-closed-form-polynomial}; see also the
Hurwitz-Bernoulli identification in Thm~\ref{thm:z-g-leading-coefficient-bernoulli}.
\end{remark}

\begin{remark}[low-genus witnesses]
\label{rem:low-genus}
Genus $g=2$: $Z_2(k) = (n/2) \sum_{j=1}^{n-1} \csc^2(\pi j/n) = (n/2)\cdot T_1(n^2) = n(n^2-1)/6$.
At $k=1, n=3$: $Z_2(1) = 3\cdot 8/6 = 4$. Genus $g=3$: $Z_3(k) = (n/2)^2 T_2(n^2) = n^2(n^2-1)(n^2+11)/180$.
At $k=1, n=3$: $Z_3(1) = 9\cdot 8\cdot 20/180 = 1440/180 = 8$. Both values match direct
counts of conformal blocks and the Verlinde fusion formula over $\widehat{\mathfrak{sl}}_2$.
\end{remark}

\section{Factored polynomial form: Theorem~\ref{thm:z-g-polynomial-form}}

\begin{theorem}[Polynomial factorisation of $Z_g$]
\label{thm:z-g-polynomial-form}
\ClaimStatusProvedHere
For $g \ge 2$ and $n = k+2$,
\begin{equation}
\label{eq:pg-factored}
Z_g(k) \;=\; P_g(n) \;=\; n^{g-1}(n^2-1)\,R_{g-2}(n^2),
\end{equation}
where $R_{g-2} \in \mathbb{Q}[x]$ has degree $g-2$ (convention: $R_{-1} = 1$ if
$g=1$; $R_0$ is a rational constant if $g=2$). The factor $n^2 - 1$ is uniform
in $g$; the residual $R_{g-2}(n^2)$ carries all genus-dependent data of degree
greater than one in $n^2$.
\end{theorem}

\begin{proof}
From Thm~\ref{thm:z-g-closed-form-polynomial},
$Z_g(k) = 2^{-(g-1)} n^{g-1} T_{g-1}(n^2)$. It suffices to show
$T_{g-1}(n^2)$ is divisible by $n^2 - 1$ with quotient a polynomial of degree
$g-2$ in $n^2$, up to the normalising $2^{-(g-1)}$.

Evaluate $T_{g-1}(n^2)$ at $n = 1$. From~\eqref{eq:zg-csc}, the cosecant sum
$\sum_{j=1}^{n-1} \csc^{2(g-1)}(\pi j / n)$ has an empty range at $n = 1$,
so $S_{g-1}(1) = 0$ and therefore $T_{g-1}(1) = 0$. Since $T_{g-1}$ is a
polynomial in $n^2$, vanishing at $n=1$ is equivalent to vanishing at
$n^2 = 1$, so $(n^2 - 1) \mid T_{g-1}(n^2)$ in $\mathbb{Q}[n^2]$.

Write $T_{g-1}(n^2) = (n^2 - 1)\,\widetilde R_{g-2}(n^2)$ with $\widetilde
R_{g-2}$ of degree $(g-1) - 1 = g-2$ in $n^2$. Absorb the prefactor
$2^{-(g-1)}$ into $R_{g-2}(n^2) := 2^{-(g-1)}\,\widetilde R_{g-2}(n^2)$; this
is again a degree-$(g-2)$ polynomial in $n^2$ with rational coefficients. Then
$Z_g(k) = n^{g-1}(n^2 - 1)\,R_{g-2}(n^2)$. \checkmark

Total degree in $n$: $(g-1) + 2 + 2(g-2) = 3g - 3$, matching
Thm~\ref{thm:z-g-closed-form-polynomial}. \checkmark
\end{proof}

\begin{remark}[explicit $R_{g-2}$ for $g \le 5$]
\label{rem:r-explicit}
Combining~\eqref{eq:T1}--\eqref{eq:T3} with the absorption $R_{g-2}(x) =
2^{-(g-1)}\,T_{g-1}(x)/(x-1)$:
\begin{align}
R_0(x) &= \tfrac{1}{6}, \\
R_1(x) &= \tfrac{1}{180}(x + 11), \\
R_2(x) &= \tfrac{1}{7560}(2 x^2 + 23 x + 191).
\end{align}
The leading coefficients $1/6, 1/180, 2/7560 = 1/3780$ match
Thm~\ref{thm:z-g-leading-coefficient-bernoulli}: for
$g=2,3,4$ the formula
$(-1)^g2^{g-1}B_{2g-2}/(2g-2)!$ gives respectively
$1/6$, $1/180$, and $1/3780$.
\end{remark}

\begin{remark}[uniform $(n^2-1)$ factor]
\label{rem:uniform-factor}
The factor $n^2 - 1$ vanishes at $n = 1$ (equivalently $k = -1$, outside the
physical range $k \ge 0$). It encodes the removal of the trivial
representation from the Verlinde sum when the level is too small to
support any nontrivial fusion; both $Z_g$ and its cosecant presentation
share this pole-at-$n=1$ structure.
\end{remark}

\section{Bernoulli leading term: Theorem~\ref{thm:z-g-leading-coefficient-bernoulli}}

\begin{theorem}[Hurwitz--Bernoulli leading coefficient]
\label{thm:z-g-leading-coefficient-bernoulli}
\ClaimStatusProvedHere
For $g \ge 2$, the leading coefficient of $R_{g-2}$ (as a polynomial of degree
$g-2$ in $n^2$) is
\begin{equation}
\label{eq:leading-coeff}
[n^{2(g-2)}]\,R_{g-2}(n^2)
\;=\;
(-1)^g\,\frac{2^{g-1}B_{2g-2}}{(2g-2)!},
\end{equation}
where $B_{2g-2}$ denotes the $(2g-2)$-th Bernoulli number in the sign
convention $B_2 = 1/6$, $B_4 = -1/30$, $B_6 = 1/42$, $B_8 = -1/30$,
$B_{10} = 5/66$, $B_{12} = -691/2730$, $B_{14} = 7/6$,
$B_{16} = -3617/510$.
\end{theorem}

\begin{proof}
Step 1 (Hurwitz cosecant series). Hurwitz (1890) proved the
$\cot$/$\csc$ power-series expansion
\begin{equation}
\label{eq:hurwitz-csc}
\left(\tfrac{\pi x}{\sin \pi x}\right)^{2}
\;=\; 1 + \sum_{m \ge 1} \frac{(2m)(2^{2m}-2)|B_{2m}|}{(2m)!}(\pi x)^{2m}
\;=\; 1 + \sum_{m\ge 1} c_m x^{2m},
\end{equation}
with
\begin{equation}
c_m \;=\; \frac{(2^{2m}-2)|B_{2m}|}{(2m-1)!}\pi^{2m}.
\end{equation}
Equivalently, for fixed $r \ge 1$,
\begin{equation}
\label{eq:csc-to-bernoulli}
\csc^{2r}(\pi x) \;=\; \frac{1}{(\pi x)^{2r}}\,\Bigl[1 + O(x^2)\Bigr]^{2r}
\;=\; \frac{1}{(\pi x)^{2r}}\,\bigl(1 + r\,c_1 x^2 + \cdots\bigr),
\end{equation}
valid as $x \to 0$.

Step 2 (leading behaviour of $S_{g-1}(n)$). Set $r = g-1$ and sum
$x = j/n$ for $j = 1, \ldots, n-1$. The leading contribution to
$S_{g-1}(n)$ as $n \to \infty$ is
\begin{equation}
\label{eq:s-asymp}
S_{g-1}(n) \;=\; \sum_{j=1}^{n-1} \csc^{2(g-1)}(\pi j/n)
\;\sim\; \frac{n^{2(g-1)}}{\pi^{2(g-1)}}\,\zeta(2g-2)\cdot 2 \quad(n \to \infty),
\end{equation}
where we used $\sum_{j=1}^{\infty} j^{-2(g-1)} = \zeta(2g-2)$ and the factor
$2$ accounts for the two tails $j \to 0^+$ and $j \to n^-$ (equivalently,
$\csc^{2(g-1)}(\pi j/n) = \csc^{2(g-1)}(\pi(n-j)/n)$, so the sum doubles each
tail).

Step 3 (Euler's even-zeta identity). Euler's formula
\begin{equation}
\zeta(2m) \;=\; \frac{(-1)^{m+1}(2\pi)^{2m} B_{2m}}{2\,(2m)!}
\;=\; \frac{(2\pi)^{2m}|B_{2m}|}{2\,(2m)!}
\end{equation}
(noting $(-1)^{m+1}B_{2m} = |B_{2m}|$ for $m \ge 1$) with $m = g-1$ gives
\begin{equation}
\zeta(2g-2) \;=\; \frac{(2\pi)^{2g-2}|B_{2g-2}|}{2\,(2g-2)!}.
\end{equation}
Substituting into~\eqref{eq:s-asymp}:
\begin{equation}
\label{eq:s-leading}
S_{g-1}(n) \;\sim\; \frac{n^{2g-2}}{\pi^{2g-2}}\cdot 2 \cdot
\frac{(2\pi)^{2g-2}|B_{2g-2}|}{2(2g-2)!}
\;=\; \frac{2^{2g-2}|B_{2g-2}|}{(2g-2)!}\,n^{2g-2}.
\end{equation}

Step 4 (leading coefficient of $Z_g$ in $n$). Combining with
Thm~\ref{thm:z-g-closed-form-polynomial}:
\begin{equation}
Z_g(k) \;=\; \left(\tfrac{n}{2}\right)^{g-1} S_{g-1}(n)
\;\sim\; \frac{1}{2^{g-1}}\cdot \frac{2^{2g-2}|B_{2g-2}|}{(2g-2)!}\,n^{3g-3}
\;=\; \frac{2^{g-1}|B_{2g-2}|}{(2g-2)!}\,n^{3g-3}.
\end{equation}

Step 5 (leading coefficient of $R_{g-2}$). From
Thm~\ref{thm:z-g-polynomial-form}, $Z_g(k) = n^{g-1}(n^2-1)\,R_{g-2}(n^2)$.
The leading $n^{3g-3}$ term comes from $n^{g-1}\cdot n^2 \cdot
[n^{2(g-2)}]\,R_{g-2}(n^2)\cdot n^{2(g-2)}$:
\begin{equation}
[n^{3g-3}]\,Z_g(k) \;=\; [n^{2(g-2)}]\,R_{g-2}(n^2).
\end{equation}
Therefore
\begin{equation}
[n^{2(g-2)}]\,R_{g-2}(n^2)
\;=\;
\frac{2^{g-1}|B_{2g-2}|}{(2g-2)!}.
\end{equation}
Since $B_{2g-2}=(-1)^g|B_{2g-2}|$, this is exactly
\eqref{eq:leading-coeff}. \checkmark
\end{proof}

\begin{remark}[low-genus verification]
\label{rem:low-genus-bernoulli}
At $g=2$, \eqref{eq:leading-coeff} gives
$2B_2/2!=1/6$, matching $R_0=1/6$. At $g=3$ it gives
$-4B_4/4!=1/180$, matching $[x]R_1(x)$. At $g=4$ it gives
$8B_6/6!=1/3780$, matching $[x^2]R_2(x)$. These three cases fix the
normalisation against the explicit polynomials.
\end{remark}

\begin{remark}[Witten symplectic volume check]
\label{rem:witten-check}
Witten's 1991 formula for the symplectic volume of $\mathcal{M}_g(\mathrm{SU}(2))$
reads $\mathrm{vol}(\mathcal{M}_g) = 2\,(2\pi^2)^{g-1}\zeta(2g-2)/(g-1)!$.
The $\zeta(2g-2)$ factor, via Euler, reintroduces $|B_{2g-2}|$, and the
large-$k$ limit of $Z_g(k)/n^{3g-3}$ matches this symplectic volume up to a
universal combinatorial normalisation. This is a second, independent
derivation of the Bernoulli leading coefficient, disjoint from the cosecant
series (step 1 above).
\end{remark}

\section{Kummer-type congruences: Theorem~\ref{thm:z-g-kummer-congruence}}

\begin{theorem}[Irregular-prime witnesses]
\label{thm:z-g-kummer-congruence}
\ClaimStatusProvedHere
Let $[n^{3g-3}]\,Z_g(k)$ denote the leading coefficient of $Z_g$ as a
polynomial in $n$; by
Theorem~\ref{thm:z-g-leading-coefficient-bernoulli} this coefficient equals
$(-1)^g2^{g-1}B_{2g-2}/(2g-2)!$. Then:
\begin{enumerate}
\item ($g=7$ witness) $B_{12} = -691/2730$, so the prime $691$ divides the
numerator of the $R_5$ leading coefficient; hence $691$ divides the
numerator of the leading coefficient of $Z_7(k)$ in $n$.
\item ($g=9$ witness) $B_{16} = -3617/510$, so the prime $3617$ divides the
numerator of the $R_7$ leading coefficient; hence $3617$ divides the
numerator of the leading coefficient of $Z_9(k)$ in $n$.
\end{enumerate}
Both primes are \emph{irregular} in the sense of Kummer: $691$ divides the
numerator of $B_{12}$, and $3617$ divides the numerator of $B_{16}$. These
divisibilities are falsification tests: if computed leading coefficients
fail to be divisible by $691$ (resp.\ $3617$), the closed-form identification
of Thms~\ref{thm:z-g-polynomial-form}--\ref{thm:z-g-leading-coefficient-bernoulli}
is wrong.
\end{theorem}

\begin{proof}
From Thm~\ref{thm:z-g-leading-coefficient-bernoulli}, the leading coefficient
of $R_{g-2}$ in $n^2$ equals
$(-1)^g2^{g-1}B_{2g-2}/(2g-2)!$.

(1) At $g = 7$: $2g-2 = 12$, $g-1 = 6$. Thus
\begin{equation}
[n^{10}]\,R_5(n^2)
\;=\;
-\,\frac{64B_{12}}{12!}
\;=\;
\frac{691}{20432412000}.
\end{equation}
The numerator is $691$; hence $691$ divides the numerator of
$[n^{18}]\,Z_7(k)$.

(2) At $g = 9$: $2g-2 = 16$, $g-1 = 8$. Thus
\begin{equation}
[n^{14}]\,R_7(n^2)
\;=\;
-\,\frac{256B_{16}}{16!}
\;=\;
\frac{3617}{41682120480000}.
\end{equation}
The numerator is $3617$; hence $3617$ divides the numerator of
$[n^{24}]\,Z_9(k)$.

Kummer's theorem identifies an odd prime $p$ as \emph{irregular} iff $p$
divides the numerator of some $B_{2m}$ with $2 \le 2m \le p-3$. For $p=691$,
$m=6$ gives $B_{12}$ with numerator $-691$; for $p=3617$, $m=8$ gives
$B_{16}$ with numerator $-3617$. These are the two earliest
irregular-prime witnesses in Bernoulli-index order for this chapter, not
the first irregular primes by size. The size order begins with
$37,59,67,\ldots$, whose Bernoulli witnesses occur at later indices. The
divisibility statements above therefore expose irregular-prime structure
directly in the Verlinde polynomial.

\emph{Falsification}. If a numerical computation of $Z_7(k)$ for a range
$k = 1, 2, \ldots, K$ yields a fitted leading coefficient whose numerator is
coprime to $691$, then one of the three theorems
Thms~\ref{thm:z-g-closed-form-polynomial}--\ref{thm:z-g-leading-coefficient-bernoulli}
is wrong. Symmetrically for $3617$ and $Z_9$. The rational identities
$B_{12}=-691/2730$ and $B_{16}=-3617/510$, together with the closed
form~\eqref{eq:pg-factored}, give the leading coefficients displayed
above.
\end{proof}

\begin{remark}[next witnesses]
\label{rem:next-witnesses}
The Bernoulli-index order and the prime-size order are different.  After
$B_{12}$ and $B_{16}$, the next low-index Bernoulli witnesses are
\[
  B_{18}=\frac{43867}{798},\qquad
  B_{20}=-\frac{283\cdot 617}{330},\qquad
  B_{22}=\frac{11\cdot 131\cdot 593}{138}.
\]
The factor $11$ in $B_{22}$ is regular, since its index violates
Kummer's bound $2m\leq p-3$; the factors $43867,283,617,131,593$ are
irregular.  By contrast, the first irregular primes by size have later
Bernoulli witnesses:
\[
  37\mid \mathrm{num}(B_{32}),\qquad
  59\mid \mathrm{num}(B_{44}),\qquad
  67\mid \mathrm{num}(B_{58}).
\]
Each such numerator prime gives a Kummer-type divisibility witness for
the leading coefficient of the corresponding $Z_g$. The sequence of
denominators of $B_{2m}$ follows the von Staudt--Clausen theorem,
$\mathrm{denom}(B_{2m}) = \prod_{(p-1) \mid 2m} p$, and this too is visible
in the denominators of $R_{g-2}$-leading coefficients.
\end{remark}

\begin{remark}[status of the falsification pipeline]
\label{rem:falsification-pipeline}
The closed form~\eqref{eq:pg-factored}, the Bernoulli leading
term~\eqref{eq:leading-coeff}, and the Kummer witnesses~$691 \mid Z_7$,
$3617 \mid Z_9$ together pin $Z_g(k)$ to a single polynomial template for
each $g$. Any deviation in any term (bulk or leading) is detectable by a
finite SymPy computation: compute $Z_g(k)$ for $k = 1, \ldots, 3g-3$ via the
Verlinde formula, interpolate as a polynomial in $n$, and verify the
factorisation, the leading coefficient, and the prime divisibility. This is
the preferred falsification protocol for this chapter.
\end{remark}

\section{Independent verification anchors}
\label{sec:zg-independent-verification}

The four theorem proofs above are checked against independent mathematical
routes.

\begin{remark}[Verification anchor: Verlinde closed form]
\label{rem:zg-verify-closed-form}
Theorem~\ref{thm:z-g-closed-form-polynomial} starts from the
Verlinde~1988 fusion formula and the modular matrix
$S_{\ell m}=\sqrt{2/n}\sin(\pi(\ell+1)(m+1)/n)$ for
$\widehat{\mathfrak{sl}}_2$. The low-genus values
$Z_2(1)=4$ and $Z_3(1)=8$ are recovered directly from the fusion ring, and
the Witten~1991 symplectic-volume asymptotic gives the same
$n^{3g-3}$ leading degree.
\end{remark}

\begin{remark}[Verification anchor: polynomial factorisation]
\label{rem:zg-verify-polynomial}
The factor $n^2-1$ in Theorem~\ref{thm:z-g-polynomial-form} follows from
the empty cosecant sum at $n=1$. The explicit polynomials
in Remarks~\ref{rem:low-genus} and~\ref{rem:r-explicit} exhibit the
factor in genera $2,3,4$, and interpolation from Verlinde values for
$k=0,\ldots,3g-3$ checks the same factorisation for $2\leq g\leq9$.
\end{remark}

\begin{remark}[Verification anchor: Bernoulli leading coefficient]
\label{rem:zg-verify-leading}
Theorem~\ref{thm:z-g-leading-coefficient-bernoulli} has two disjoint
derivations: Hurwitz's cosecant expansion together with Euler's
even-zeta identity, and Witten's large-level symplectic-volume formula.
Both give
\[
  [n^{2(g-2)}]R_{g-2}(n^2)
  =
  (-1)^g\frac{2^{g-1}B_{2g-2}}{(2g-2)!}.
\]
The explicit values $1/6$, $1/180$, and $1/3780$ at $g=2,3,4$ verify
the normalisation against the displayed polynomials.
\end{remark}

\begin{remark}[Verification anchor: Kummer witnesses]
\label{rem:zg-verify-kummer}
Theorem~\ref{thm:z-g-kummer-congruence} uses
$B_{12}=-691/2730$ and $B_{16}=-3617/510$. Direct rational reduction gives
\[
  [n^{18}]Z_7(k)=\frac{691}{20432412000},
  \qquad
  [n^{24}]Z_9(k)=\frac{3617}{41682120480000}.
\]
The prime divisibilities are therefore visible in the reduced leading
coefficients themselves.
\end{remark}

\section{Dependencies}

\begin{itemize}
\item Theorem~\ref{thm:z-g-closed-form-polynomial} uses the
Verlinde~1988 formula; Witten's 1991 volume asymptotic and the
low-genus fusion-ring counts check the normalisation.
\item Theorem~\ref{thm:z-g-polynomial-form} uses the empty cosecant sum
at $n=1$ and the polynomial structure of cosecant power sums; the
low-genus polynomials check the factor $n^2-1$.
\item Theorem~\ref{thm:z-g-leading-coefficient-bernoulli} uses
Hurwitz's 1890 cosecant series and Euler's even-zeta identity; Witten's
large-level volume formula gives the same Bernoulli factor.
\item Theorem~\ref{thm:z-g-kummer-congruence} uses von Staudt--Clausen
and Kummer irregularity; direct interpolation of the Verlinde polynomial
checks the displayed leading numerators.
\end{itemize}

This chapter's closed form is the $G = \mathrm{SL}_2$
specialisation of the universal Verlinde polynomial $P_g$. The
Hurwitz--Bernoulli leading coefficient generalises to arbitrary simple
$G$ via the Macdonald denominator formula and the Kac--Peterson
character-theoretic derivation of $Z_g^G(k)$; this chapter fixes only the
$\mathrm{SL}_2$ normalisation.

\section{Arithmetic duality with the Virasoro shadow tower}
\label{sec:zg-arithmetic-duality}

The four theorems of this chapter pin the arithmetic of $Z_g(k)$ on
$\mathrm{Bun}_{\mathrm{SL}_2}(\Sigma_g)$ through the Hurwitz--Bernoulli
leading coefficient
$(-1)^g2^{g-1}B_{2g-2}/(2g-2)!$ and its Kummer-irregular prime
witnesses $\{691, 3617, \ldots\}$. The Virasoro shadow tower
$\{S_r(\Vir_c)\}_{r \geq 2}$ carries the complementary arithmetic:
the two invariants are geometrically disjoint (different localisations),
and the two
\emph{Bernoulli-leading} Kummer witnesses $\{691, 3617\}$ that appear
on the $Z_g$ side at genera $g = 7, 9$ are provably \emph{absent} on
the $S_r$ side through $r = 11$. Kummer-irregular primes ordered by size
belong to a different filtration; for instance $2111$ can appear in an
$S_r$ coefficient without contradicting the stated Bernoulli-index
duality.

\begin{theorem}[$Z_g$ vs $S_r(\Vir_c)$ arithmetic duality at the
Bernoulli-leading Kummer pair]
\label{thm:z-g-s-r-arithmetic-duality}
\ClaimStatusProvedHere
Let
\[
  \mathrm{IrrKum}_{\mathrm{Bern},2}:=\{691,3617\}
\]
denote the Kummer-irregular primes appearing at the two smallest
Bernoulli indices that contain irregular numerator factors:
$691\mid\mathrm{num}(B_{12})$ and
$3617\mid\mathrm{num}(B_{16})$. Then:
\begin{enumerate}
\item (Presence on the $Z_g$ side.) The prime $691$ divides the reduced
integer numerator of the leading coefficient $[n^{18}]Z_7(k)$, and the
prime $3617$ divides the reduced integer numerator of
$[n^{24}]Z_9(k)$. Equivalently,
\[
  691\mid \mathrm{num}(B_{12}),\qquad
  3617\mid \mathrm{num}(B_{16}),
\]
and in both cases the scalar multiplier
$(-1)^g2^{g-1}/(2g-2)!$ is a $p$-adic unit.
\item (Absence on the $S_r$ side through $r = 11$.) For every $r$ in the
range $2 \leq r \leq 11$, neither $691$ nor $3617$ divides the numerator
of $S_r(\Vir_c) \cdot c^{r-3} (5c + 22)^{\lfloor (r-2)/2 \rfloor}$
regarded as an integer polynomial in $c$. Explicitly
\[
  691,\, 3617 \;\nmid\; \mathrm{num}\bigl[\,S_r(\Vir_c) \cdot c^{r-3}
  (5c+22)^{\lfloor (r-2)/2 \rfloor}\,\bigr]
  \qquad\text{for all } 2 \leq r \leq 11.
\]
\end{enumerate}
In particular the two Bernoulli-leading Kummer witnesses $691$ (seen at
$g = 7$, $B_{12}$) and $3617$ (seen at $g = 9$, $B_{16}$) are
\emph{present} in the $Z_g$ leading coefficients at those genera and
\emph{absent} from every $S_r(\Vir_c)$ coefficient through weight
$r = 11$. The duality is sharp for this Bernoulli-index pair; it is not
a global absence theorem for Kummer-irregular primes in the Virasoro
shadow tower.
\end{theorem}

\begin{proof}
Part (1) is a direct composition of
Theorem~\ref{thm:z-g-leading-coefficient-bernoulli}
(leading coefficient of $R_{g-2}$ equals
$(-1)^g2^{g-1}B_{2g-2}/(2g-2)!$) with the Kummer-irregularity criterion
for the pair $\{691, 3617\}$:
$691 \mid \mathrm{num}(B_{12})$ (with $g = 7$, $2g - 2 = 12$) and
$3617 \mid \mathrm{num}(B_{16})$ (with $g = 9$, $2g - 2 = 16$).
Explicit witnesses at $g = 7$ and $g = 9$ are
Theorem~\ref{thm:z-g-kummer-congruence} (1) and (2). The factor
$2^{g-1}/(2g-2)!$ has no factor $691$ at $g=7$ and no factor $3617$ at
$g=9$, so the Bernoulli numerator survives rational reduction.

Part (2) is exactly Theorem~\ref{thm:s-r-kummer-absent-through-r-11},
which pins the
$\mathbb{Z}$-polynomial numerator
$N_r(c) = S_r(\Vir_c) \cdot c^{r-3} (5c+22)^{\lfloor (r-2)/2 \rfloor}$
through $r = 11$ and confirms that $691, 3617$ do not appear among the
primes in $\mathcal{P}_{\le 11}$. Exact integer factorisation of every
coefficient of $N_2,\ldots,N_{11}$ checks the claim.

The two statements together constitute the duality at the two
Bernoulli-leading Kummer witnesses: $\{691, 3617\}$ are present on the
$Z_g$ side (via the Hurwitz--Bernoulli leading factor) and absent on the $S_r$
side (through $r = 11$, verified). The statement does not extend to
every Kummer-irregular prime: the prime $2111$ is Kummer-irregular (via
$2111 \mid \mathrm{num}(B_{1038})$) and simultaneously
$2111 \mid 29554$, with $29554$ the constant-term coefficient of $N_9$.
This coincidence does not violate the stated duality because $2111$ is
not witnessed on the $Z_g$ side: the $g$-th Kummer-presence is
controlled by $\mathrm{num}(B_{2g-2})$, and for $g \leq 9$ the relevant
Bernoulli numerators $\mathrm{num}(B_{2}), \ldots, \mathrm{num}(B_{16})$
have only $691$ and $3617$ as Kummer-irregular prime factors.
\checkmark
\end{proof}

\begin{remark}[geometric disjointness of localisations]
\label{rem:localization-disjointness}
Theorem~\ref{thm:z-g-s-r-arithmetic-duality} expresses a
\emph{geometric} disjointness underlying the arithmetic duality. The two
sides localise on different spaces.

\begin{itemize}
\item The $Z_g(k)$ leading coefficient is a characteristic number on
$\mathrm{Bun}_{\mathrm{SL}_2}(\Sigma_g)$, the moduli stack of
$\mathrm{SL}_2$-bundles on the genus-$g$ surface $\Sigma_g$. By
Theorem~\ref{thm:z-g-leading-coefficient-bernoulli} together with
Remark~\ref{rem:witten-check}, the leading coefficient is the
large-$k$ symplectic volume of the moduli space of flat
$\mathrm{SU}(2)$-connections, computed by Witten~1991 via a ratio of
zeta values; the Bernoulli content enters through Euler's even-zeta
identity. The arithmetic of $Z_g$ is thus entirely a moduli-of-bundles
invariant and its primes are controlled by the $B_{2m}$ denominator
spectrum (von Staudt--Clausen) and numerator spectrum (Kummer
irregularity).
\item The Virasoro shadow coefficient $S_r(\Vir_c)$ is a
Chiral Hochschild residue of the chiral algebra itself. By
Theorem~\ref{thm:virasoro-shadow-recurrence}, $S_r$ is a
Riccati-propagated rational function of $c$ starting from the initial
data $(S_2, S_3, S_4, S_5)$, with no input from any moduli stack.
$S_r(\Vir_c)$ localises on the universal Maurer--Cartan element
$\Theta_{\Vir_c}$ in the ordered convolution dGLA -- inherently a
statement on $\ChirHoch(\Vir_c)$, not on
$\mathrm{Bun}_G(\Sigma_g)$.
\end{itemize}

The two sides never share a common geometric localisation. In particular
there is no universal bridge that would force the Kummer-irregular primes
of $Z_g$ to appear in $S_r$ or vice versa. The arithmetic duality
(Kummer-present on $Z_g$; Kummer-absent on $S_r$) is the shadow of this
geometric disjointness at the arithmetic level.

This disjointness resolves a potentially misleading reading of the
shadow-Feynman bridge $F_g \leftrightarrow S_{2g-2}$: the bridge is a
\emph{structural} identification of the all-genus free energy with the
shadow tower under the scalar-primary channel, not an equality of
\emph{arithmetics}. The Bernoulli content of $F_g$ (via $B_{2g-2}$)
enters through the genus normalisation, which is a moduli-of-bundles
feature; the Riccati content of $S_r$ enters through the ordered
MC recursion, which is a chiral-algebra feature. The two converge on a
single formal power series under the bridge, but the arithmetic of their
coefficients is distinct.
\end{remark}

\begin{remark}[status of extensions]
\label{rem:zg-sr-duality-extensions}
The theorem pins the duality for the Bernoulli-index pair
$B_{12},B_{16}$ and for the shadow range $r \leq 11$.  It does not
assert absence of all Kummer-irregular primes from the Virasoro tower:
$2111$ already appears in $N_9$.  The higher Bernoulli-index witnesses
belong to Theorem~\ref{thm:higher-kummer-refined-duality}; the Laurent
tier refinement belongs to Theorem~\ref{thm:kummer-laurent-depth-controlled}.
\end{remark}

\begin{remark}[Verification anchor: arithmetic duality]
\label{rem:zg-verify-sr-duality}
Theorem~\ref{thm:z-g-s-r-arithmetic-duality} combines two independent
computations. On the $Z_g$ side, the reduced leading coefficients at
$g=7,9$ contain the Bernoulli numerators $691$ and $3617$. On the
Virasoro side, Theorem~\ref{thm:s-r-kummer-absent-through-r-11} factors
every coefficient of
$N_r(c)=S_r(\Vir_c)c^{r-3}(5c+22)^{\lfloor(r-2)/2\rfloor}$ for
$2\leq r\leq11$ and finds neither prime. The first computation is a
moduli-of-bundles calculation; the second is the Riccati transport in the
ordered chiral Hochschild recursion.
\end{remark}

\bigskip

\emph{End of chapter.}
